\documentclass[12pt]{article}

% --- Preámbulo (Configuración) ---
\usepackage[utf8]{inputenc}       % Permite usar tildes y eñes
\usepackage[spanish]{babel}       % Configura el idioma a español
\usepackage{amsmath, amssymb}     % Paquetes para símbolos matemáticos
\usepackage{graphicx}             % Para insertar imágenes

\title{Mi Primer Documento en \LaTeX}
\author{Tu Nombre}
\date{\today}

% --- Inicio del Documento ---
\begin{document}

\maketitle % Imprime el título, autor y fecha

\tableofcontents % Genera el índice automático
\newpage

\section{Introducción}
Este es un ejemplo básico de un documento creado en \LaTeX. A diferencia de los editores de texto convencionales, aquí escribimos código que luego se compila para generar un PDF profesional.

\section{Conceptos Básicos}
En esta sección mostramos algunas funcionalidades útiles:

\subsection{Listas}
Podemos crear listas fácilmente:
\begin{itemize}
    \item Primero, aprender la sintaxis.
    \item Segundo, practicar la compilación.
    \item Tercero, organizar el contenido.
\end{itemize}

\subsection{Matemáticas}
Una de las mayores potencias de \LaTeX{} es la escritura de fórmulas. Podemos poner ecuaciones en línea como $E = mc^2$ o centradas en su propia línea:

\[
\int_{a}^{b} x^2 \, dx = \frac{b^3 - a^3}{3}
\]

\section{Conclusión}
\LaTeX{} es ideal para documentos académicos, tesis y artículos científicos debido a su excelente gestión de bibliografía y formato.

\end{document}